% P8 hybrid architecture: sensor + scorer + scribe, post-score agentic triage
\section{The Hybrid Architecture: Sensor and Scribe, Not Scorer}
\label{sec:p8-architecture}

The taxonomy dictates the architecture. Each model family holds exactly the
seats it wins, and the seams between them are typed interfaces rather than
shared prompts (\figref{fig:p8-architecture}).

\begin{figure}[t]
\centering
\resizebox{\linewidth}{!}{%
\begin{tikzpicture}[
  font=\small,
  node distance=7mm and 9mm,
  box/.style={draw, rounded corners=2pt, align=center, minimum height=9mm,
              inner sep=4pt, text width=30mm},
  llm/.style={box, fill=blue!8},
  tree/.style={box, fill=green!10},
  human/.style={box, fill=orange!12},
  data/.style={box, fill=gray!8, text width=26mm},
  arr/.style={-{Latex[length=2mm]}, thick}
]
% Inputs
\node[data] (txn) {Transaction stream\\(fixed schema)};
\node[data, below=of txn] (docs) {Documents / images\\(disputes, onboarding)};
% Sensor
\node[llm, right=of docs] (sensor) {\textbf{LLM sensor}\\VLM feature extraction\\$\to$ typed features};
% Scorer
\node[tree, right=of txn, xshift=42mm] (scorer) {\textbf{XGBoost scorer}\\real-time, calibrated\\fraud score};
% Alert queue + triage
\node[llm, below=of scorer] (triage) {\textbf{LLM agentic triage}\\post-score alert-queue\\investigation};
% Scribe
\node[llm, right=of triage, xshift=2mm] (scribe) {\textbf{LLM scribe}\\SAR narrative /\\adverse-action draft};
% Human
\node[human, above=of scribe] (analyst) {\textbf{Human analyst}\\review, decide, file};
% Edges
\draw[arr] (txn) -- (scorer);
\draw[arr] (docs) -- (sensor);
\draw[arr] (sensor.east) -- node[above, sloped, font=\scriptsize]{numeric features} (scorer.west |- sensor.east) -- (scorer.south west);
\draw[arr] (scorer) -- node[right, font=\scriptsize]{alerts} (triage);
\draw[arr] (triage) -- node[below, font=\scriptsize]{case file} (scribe);
\draw[arr] (scribe) -- node[right, font=\scriptsize]{draft} (analyst);
\draw[arr] (analyst.west) to[out=180, in=45, looseness=1.1]
  node[above, sloped, font=\scriptsize]{labels} (scorer.north east);
\end{tikzpicture}%
}
\caption{The hybrid architecture. The LLM never sits in the synchronous scoring
path: the \emph{sensor} converts unstructured evidence to typed numeric
features upstream of the tree; the calibrated XGBoost \emph{scorer} owns the
real-time decision; the \emph{agentic triage} loop and the \emph{scribe}
operate post-score on the alert queue; the human analyst reviews drafts,
decides, files, and returns confirmed labels to retrain the scorer.}
\label{fig:p8-architecture}
\end{figure}

\paragraph{LLM as sensor.} Upstream of the tree, a vision--language model
consumes the unstructured evidence of \secref{sec:p8-gap-docs} and emits
\emph{typed features}: document-consistency scores, extracted field values,
tamper-likelihood indicators, mismatch flags. The interface is a fixed numeric
schema, so from the tree's perspective the VLM is just another feature source.
This placement has two consequences. First, the tree's accuracy and
calibration properties are preserved: the score is still produced by the model
family that wins the scorer seat. Second, the injection surface is contained:
adversarial text inside a document can at worst corrupt that document's
feature values---a bounded, per-case failure the tree averages over---rather
than hijack a decision-making agent, since the sensor has no tools, no memory
across cases, and no authority \citep{greshake2023injection}.

\paragraph{XGBoost as scorer.} The synchronous authorization path is unchanged
from the incumbent design: fixed-schema features (now augmented by sensor
outputs where evidence exists) enter a boosted tree; a calibrated score exits
in milliseconds at near-zero marginal cost. Nothing in the real-time loop
parses natural language. All four operational arguments of
\secref{sec:p8-scorer} are honored by construction.

\paragraph{Post-score agentic triage.} Alerts crossing the threshold enter a
queue where an LLM agent performs the mechanical investigation of
\secref{sec:p8-gap-agentic}: retrieving history, cross-referencing entities,
checking narrative consistency, and ranking the queue by investigative
priority with a written rationale per case. The agent is asynchronous
(seconds-to-minutes latency is acceptable), budget-capped (the token spend
lands only on the $\sim$1\% of traffic that alerts), and advisory (it cannot
close a case or file a report). Its cold-start mode (\secref{sec:p8-gap-coldstart})
runs the same loop seeded by typology bulletins instead of scorer alerts,
covering the label vacuum until confirmed cases retrain the tree.

\paragraph{LLM as scribe.} At the end of the queue, the scribe drafts the two
regulated documents of \secref{sec:p8-gap-narration}: the SAR narrative,
conditioned on the case file and structured to FinCEN's
who/what/when/where/why guidance \citep{fincen2003sarnarrative}; and, in
credit-decision contexts, the adverse-action statement of specific reasons,
conditioned on the tree's feature attributions so that the stated reasons are
the model's actual principal reasons as ECOA/Regulation~B demands
\citep{cfr1002_9, cfpb2022circular}. Every draft carries provenance links from
each claim to its evidence, and a human analyst signs before anything is
filed. The scribe is the highest-value, lowest-risk seat: pure text
generation, fully reviewable, in a category with no non-LLM competitor.

\paragraph{The label loop.} Analyst dispositions flow back as confirmed
labels, retraining the scorer and grading the sensor and triage components.
The architecture is thus not a static ensemble but a division of labor with a
single supervised spine: the tree learns from every closed case, and the LLM
seats are evaluated by how much they improve the tree's inputs, the queue's
ordering, and the filing's speed---each independently measurable, which is the
program discipline (``measure capability, don't trust the label'') applied to
the system's own components.
